% ============================================================
%  Chapter 1 — What Is a Smartphone?
%  Inside a Smartphone: The Tiny World in Your Hand
% ============================================================

\chapter{What Is a Smartphone?}

\chapteropening{%
  You hold it in your hand every day.
  But have you ever wondered what is going on inside?
}
\chapterbadge{Mission Start: Shrink Down and Step Inside!}

% --- Opening story ---
\section*{Dot Returns}

Sam was sitting on the steps outside, holding a smartphone.

The screen glowed. Music was playing. A photo notification popped up.
Then a map. Then a game.

Sam frowned and said:

\begin{center}
  \Large\bfseries\color{TechPurple}
  ``This little thing does so much. How does it all fit inside?''
\end{center}

\vspace{8pt}

A tiny glowing dot zoomed out from the corner of the screen with a familiar grin.

\character{Dot}{Remember me? Last time we explored the Internet together.
  Now there is a whole new adventure waiting — \emph{inside} this very phone.
  Ready? Follow me in!}

\sectionrule

% --- More than a phone ---
\section*{More Than Just a Phone}

The word ``smartphone'' has the word \emph{phone} in it — but a smartphone is so much more.

\begin{picturethis}
  Imagine a Swiss Army knife.
  It looks like a simple knife, but unfold it and you find scissors, a screwdriver,
  a saw, a magnifying glass, and more.
  A smartphone is like that — it folds dozens of devices into one thin rectangle.
\end{picturethis}

Inside one smartphone you will find:
\begin{itemize}
  \item A \term{telephone} (calling and texting)
  \item A \term{camera} (photos and videos)
  \item A \term{music player}
  \item A \term{navigation device} (maps and GPS)
  \item A \term{web browser} (for exploring websites on the internet)
  \item A \term{calculator}, \term{torch}, \term{calendar}, and much more
\end{itemize}

\begin{funfact}
  The first iPhone was released in 2007 — about 20 years ago.
  Before that, phones could only call and send text messages.
\end{funfact}

\sectionrule

% --- A tiny computer ---
\section*{A Computer That Fits in Your Pocket}

Here is the most surprising fact: a smartphone is a \term{computer}.
Not just similar to a computer — it \emph{is} one.
Any computer does three big jobs: it takes in information, works on it,
and gives you a result. Your smartphone does all three, all day long.

Like a laptop or desktop, a smartphone has:
\begin{itemize}
  \item A \term{processor} (the brain that does all the thinking)
  \item \term{Memory} (for working on things right now)
  \item \term{Storage} (for keeping things long term)
  \item \term{Input} (touchscreen, microphone, camera)
  \item \term{Output} (screen, speaker, vibration)
\end{itemize}

\begin{quickmap}
  \textbf{How powerful is a smartphone?}
  \begin{enumerate}[label=\textbf{\arabic*.}]
    \item A 1990s desktop computer could do about 100 million calculations per second.
    \item A modern smartphone processor can do over 15 \emph{billion} steps per second.
    \item That is more than 150 times faster — and it fits in your pocket.
  \end{enumerate}
\end{quickmap}

\sectionrule

% --- How smartphones changed life ---
\section*{How Smartphones Changed Everyday Life}

Before smartphones, if you wanted to:
\begin{itemize}
  \item Check the weather — you watched TV or bought a newspaper.
  \item Take a photo — you used a separate camera with film.
  \item Find directions — you read a paper map.
  \item Listen to music — you carried a separate music player.
\end{itemize}

Today, one device in your pocket does all of these things instantly.

\begin{diagram}[The many tools packed into one smartphone.]
  \begin{tikzpicture}[>=Stealth]
    % Central phone
    \node[draw=TechPurple,fill=TechPurple!15,rounded corners=8pt,
          minimum width=2.2cm,minimum height=3.8cm,line width=1.5pt] (phone) at (0,0) {};
    \node[align=center,font=\bfseries\small\color{TechPurple}] at (0,0) {Smart\\phone};
    % Surrounding labels
    \foreach \angle/\label in {
      90/Camera,
      40/GPS,
      0/Music,
      320/Browser,
      270/Calculator,
      220/Torch,
      180/Calendar,
      140/Phone
    }{
      \node[font=\small,draw=TechPurple!40,fill=TechPurple!5,
            rounded corners=3pt,inner sep=3pt]
        at (\angle:3.2cm) {\label};
      \draw[TechPurple!50,->] (\angle:2.4cm) -- (\angle:1.2cm);
    }
  \end{tikzpicture}
\end{diagram}

\sectionrule

% --- New words ---
\begin{newwords}
  \begin{description}[labelwidth=3.5cm, leftmargin=3.7cm]
    \item[\term{Smartphone}]   A mobile phone that is also a computer, camera, map, and much more.
    \item[\term{Processor}]    The brain of a device — carries out instructions. (We meet it in Chapter 2.)
    \item[\term{Memory}]       Holds information the device is working with right now.
    \item[\term{Storage}]      Keeps photos, apps, and files even when the phone is off.
    \item[\term{Input}]        The ways information goes \emph{into} a phone — touch, voice, camera.
    \item[\term{Output}]       The ways information comes \emph{out} — screen, sound, vibration.
  \end{description}
\end{newwords}

\sectionrule

% --- Try It ---
\begin{tryit}
  \textbf{Count the devices in your phone!}

  Look at the home screen of a smartphone (with a grown-up nearby).
  Find an app for each of these:
  \begin{itemize}
    \item A clock or alarm
    \item A camera
    \item A weather app
    \item A music or video player
    \item A map
  \end{itemize}
  How many did you find? Could you find even more?
\end{tryit}

\sectionrule

% --- Big Idea ---
\begin{bigidea}
  A smartphone is a \textbf{tiny computer} that combines dozens of devices
  into one. Inside, it has a brain, memory, storage, cameras, and sensors
  all working together.
\end{bigidea}

\character{Dot}{Now that you know what a smartphone is, let's meet the
  most important part inside it — the brain. Follow me to Chapter 2!}
